\documentclass[11pt]{article}

\usepackage[margin=1in]{geometry}
\usepackage{amsmath,amssymb,mathtools}
\usepackage{algorithm}
\usepackage{algpseudocode}
\usepackage{booktabs}
\usepackage[hidelinks]{hyperref}

\newcommand{\R}{\mathbb{R}}
\newcommand{\argmin}{\operatorname*{arg\,min}}
\newcommand{\dist}{\operatorname{dist}}
\newcommand{\proj}{\operatorname{Proj}}
\newcommand{\clip}{\operatorname{clip}}
\newcommand{\norm}[1]{\left\lVert #1 \right\rVert}
\newcommand{\ip}[2]{\left\langle #1,#2 \right\rangle}
\newcommand{\Lagr}{\mathcal{L}}
\newcommand{\Q}{\mathcal{Q}}

\title{Nonconvex Online ADMM for Quantized Training and LLM-Style PTQ}
\author{}
\date{}

\begin{document}
\maketitle

\section{Purpose and Positioning}

The convex ADMM analysis should not be transferred to neural-network training
by replacing the word ``convex'' with ``nonconvex.''  The right theorem target
changes.  In convex ADMM, one can compare to a global saddle point and prove
primal-dual residual decay.  In nonconvex training and quantization, the
appropriate targets are stationarity of a merit function, KKT residuals,
convergence to a critical point under a KL assumption, or convergence to a
quantization/noise floor.

The reusable part of the online-penalty proof is the bounded path length of
\(u_k=\log\rho_k\).  If \(u_k\in[u_{\min},u_{\max}]\) and the update direction
is clipped, then
\begin{equation}
    u_{k+1}
    =
    \Pi_{[u_{\min},u_{\max}]}
    (u_k-\eta_k\widetilde g_k),
    \qquad
    \norm{\widetilde g_k}\le G_u
\end{equation}
implies
\begin{equation}
    |u_{k+1}-u_k|\le \eta_k G_u.
    \label{eq:path-length}
\end{equation}
This path-length control does not require convexity of the primal problem or
convexity of the scalar meta-loss.  It only requires projection, clipping, and
compactness of \(U=[u_{\min},u_{\max}]\).  The nonconvex proof should therefore
separate the penalty-path perturbation layer from the nonconvex ADMM descent
layer.

\section{Generic Nonconvex Variable-Penalty ADMM}

Consider
\begin{equation}
    \min_{x,z}\; f(x)+g(z)
    \quad \text{s.t.}\quad
    Ax+Bz=c,
    \label{eq:generic-nonconvex}
\end{equation}
where \(f\) may be smooth nonconvex and \(g\) may be nonsmooth, weakly convex,
prox-regular, or an indicator of a discrete quantization set.  The unscaled
augmented Lagrangian is
\begin{equation}
    \Lagr_{\rho}(x,z,\lambda)
    =
    f(x)+g(z)+\ip{\lambda}{Ax+Bz-c}
    +\frac{\rho}{2}\norm{Ax+Bz-c}^2.
\end{equation}
A proof-friendly nonconvex ADMM scheme adds proximal regularization:
\begin{align}
    x^{k+1}
    &\approx
    \argmin_x
    \Lagr_{\rho_k}(x,z^k,\lambda^k)
    +\frac{\tau_x}{2}\norm{x-x^k}^2,\\
    z^{k+1}
    &\approx
    \argmin_z
    \Lagr_{\rho_k}(x^{k+1},z,\lambda^k)
    +\frac{\tau_z}{2}\norm{z-z^k}^2,\\
    \lambda^{k+1}
    &=
    \lambda^k+\rho_k(Ax^{k+1}+Bz^{k+1}-c).
\end{align}
In scaled form, \(w^k=\lambda^k/\rho_k\), changing \(\rho\) still requires
\begin{equation}
    w^{k+1}\leftarrow \frac{\rho_k}{\rho_{k+1}}w^{k+1}.
    \label{eq:nonconvex-rescale}
\end{equation}

\begin{algorithm}[t]
\caption{Nonconvex online ADMM with bounded log-penalty path}
\label{alg:nonconvex-online-admm}
\begin{algorithmic}[1]
\Require \(x^0,z^0,w^0,\rho_0\), proximal parameters \(\tau_x,\tau_z\)
\For{\(k=0,1,\ldots,K-1\)}
    \State Approximately minimize the \(x\)-subproblem with proximal term
    \State Approximately minimize or project the \(z\)-subproblem
    \State \(w^{k+1}\gets w^k+Ax^{k+1}+Bz^{k+1}-c\)
    \State Compute \(r^{k+1}=Ax^{k+1}+Bz^{k+1}-c\)
    \State Compute \(d^{k+1}=B(z^{k+1}-z^k)\)
    \State \(g_k\gets u_k+\log(\norm{d^{k+1}}+\epsilon)
           -\log(\norm{r^{k+1}}+\epsilon)\)
    \State \(u_{k+1}\gets
           \Pi_U(u_k-\eta_k\clip(g_k;G_u))\), \(\rho_{k+1}=e^{u_{k+1}}\)
    \If{\(\rho_{k+1}\ne\rho_k\)}
        \State Rescale \(w^{k+1}\gets(\rho_k/\rho_{k+1})w^{k+1}\)
    \EndIf
\EndFor
\end{algorithmic}
\end{algorithm}

\section{Task-Aware Online Losses for Nonconvex Quantization}

The residual-balancing loss is useful because it is analytic in
\(u=\log\rho\), but it is not necessarily aligned with deploy quality.  For
quantization, the observed task metric
\[
    m_k =
    \begin{cases}
        \text{deploy loss}, & \text{quantized training},\\
        \text{calibration relative error}, & \text{PTQ},
    \end{cases}
\]
can be fed back into the penalty update.  The experiments implement three
online losses.  The baseline residual loss is
\begin{equation}
    \ell_k^{\mathrm{res}}(u)
    =
    \frac{1}{2}
    \left(
        u+\log(\norm{d^{k+1}}+\epsilon)
        -\log(\norm{r^{k+1}}+\epsilon)
    \right)^2,
    \label{eq:res-online-loss-nonconvex}
\end{equation}
where \(r^{k+1}=W^{k+1}-Z^{k+1}\) and
\(d^{k+1}=Z^{k+1}-Z^k\) in the quantization split.  Its gradient is
\begin{equation}
    g_k^{\mathrm{res}}
    =
    u_k+\log(\norm{d^{k+1}}+\epsilon)
    -\log(\norm{r^{k+1}}+\epsilon).
\end{equation}

\subsection{Composite Task-Secant Loss}

Because the deploy metric is a black-box function of the previous penalty
trajectory, the code uses a one-dimensional secant estimate:
\begin{equation}
    g_k^{\mathrm{task}}
    =
    \clip\left(
    \frac{\log(m_k+\epsilon)-\log(m_{k-1}+\epsilon)}
         {u_k-u_{k-1}},
    G_{\mathrm{task}}
    \right)
    \quad \text{when } |u_k-u_{k-1}|>\epsilon,
    \label{eq:task-secant-gradient}
\end{equation}
and \(g_k^{\mathrm{task}}=0\) otherwise.  The composite update is
\begin{equation}
    u_{k+1}
    =
    \Pi_U
    \left(
        u_k-\eta_k
        \clip\left(
            \alpha g_k^{\mathrm{res}}
            +\beta g_k^{\mathrm{task}},
            G_u
        \right)
    \right).
    \label{eq:task-aware-update}
\end{equation}
This is the method labeled \texttt{online\_ogd\_task\_aware}.

\subsection{Feasibility-Guarded Task Loss}

The feasibility-guarded variant adds pressure to increase \(\rho\) when the
current primal residual is large relative to the ADMM stopping threshold
\(\varepsilon_k^{\mathrm{pri}}\).  Define
\begin{equation}
    s_k^{\mathrm{feas}}
    =
    \clip\left(
        \log
        \frac{\norm{r^{k+1}}+\epsilon}
             {\varepsilon_k^{\mathrm{pri}}+\epsilon},
        G_u
    \right).
\end{equation}
The update direction becomes
\begin{equation}
    g_k^{\mathrm{ft}}
    =
    \clip\left(
        \alpha g_k^{\mathrm{res}}
        +\beta g_k^{\mathrm{task}}
        -\chi s_k^{\mathrm{feas}},
        G_u
    \right).
    \label{eq:feas-task-update}
\end{equation}
The minus sign is intentional: if \(s_k^{\mathrm{feas}}>0\), the projected OGD
step \(u_{k+1}=u_k-\eta_k g_k^{\mathrm{ft}}\) moves \(u\) upward.  This is the
method labeled \texttt{online\_ogd\_task\_feasibility}.

\subsection{Accept/Reject Task Safeguard}

The accept/reject ablation shrinks the penalty when the task metric increases
too quickly:
\begin{equation}
    \log(m_k+\epsilon)-\log(m_{k-1}+\epsilon)>\Delta_{\max}
    \quad\Longrightarrow\quad
    u_{k+1}=\Pi_U(u_k+\log\theta),
    \qquad 0<\theta<1.
    \label{eq:task-accept-reject}
\end{equation}
Otherwise it uses the composite task-secant update
\eqref{eq:task-aware-update}.  The current experiment shows that this
one-sided shrink rule can drive \(\rho\) too low in PTQ, so it is best treated
as a diagnostic ablation rather than a final algorithm.

\begin{algorithm}[t]
\caption{Task-aware online penalty update}
\label{alg:task-aware-penalty}
\begin{algorithmic}[1]
\Require Current \(u_k\), residuals \(r^{k+1},d^{k+1}\), task metric \(m_k\)
\State Compute \(g_k^{\mathrm{res}}\) from \eqref{eq:res-online-loss-nonconvex}
\State Estimate \(g_k^{\mathrm{task}}\) from \eqref{eq:task-secant-gradient}
\If{using feasibility guard}
    \State Add the feasibility term in \eqref{eq:feas-task-update}
\ElsIf{using accept/reject and task metric increased too quickly}
    \State Shrink \(u_k\) according to \eqref{eq:task-accept-reject}
\Else
    \State Apply the composite update \eqref{eq:task-aware-update}
\EndIf
\State Return \(\rho_{k+1}=\exp(u_{k+1})\) and rescale the scaled dual
\end{algorithmic}
\end{algorithm}

\section{Robustness to Initial Penalty}

The online update should be tested over a grid of initial penalties
\[
    \rho_0\in\{10^{-2},10^{-1},1,10,10^2\}.
\]
This is more informative than reporting a single tuned start.  A projected
online method is robust in the penalty path if different initializations lead
to similar terminal values of \(u_k=\log\rho_k\).  However, robustness in
\(\rho_k\) is not the same as robustness in deploy quality.  In the current
experiments, residual-only OGD often returns to a similar final penalty from
different starts, but quantized deploy loss can still retain finite-horizon
memory of the early ADMM trajectory.

For a finite horizon \(K\), the initialization sensitivity can be summarized by
\begin{equation}
    \operatorname{Range}_\rho(\mathcal A)
    =
    \max_{\rho_0\in\mathcal G}
    M_K(\mathcal A,\rho_0)
    -
    \min_{\rho_0\in\mathcal G}
    M_K(\mathcal A,\rho_0),
    \label{eq:rho-init-range}
\end{equation}
where \(M_K\) is the terminal deploy metric and \(\mathcal G\) is the initial
penalty grid.  A good online tuner should have both a low mean deploy metric
and a small range in \eqref{eq:rho-init-range}.

\section{Vector Penalties}

For multi-block quantization, a scalar penalty may be too restrictive because
different layers have different curvature, activation scales, and
quantization sensitivity.  The vector-rho extension replaces
\eqref{eq:quant-split} by
\begin{equation}
    \min_{\{W_j,Z_j\}}\;
    \sum_{j=1}^m F_j(W_j)+I_{\Q_j}(Z_j)
    \quad\text{s.t.}\quad W_j-Z_j=0,\quad j=1,\ldots,m,
\end{equation}
with blockwise scaled ADMM updates
\begin{align}
    W_j^{k+1}
    &\approx
    \argmin_{W_j}
    F_j(W_j)
    +\frac{\rho_{j,k}}{2}
    \norm{W_j-Z_j^k+U_j^k}^2,\\
    Z_j^{k+1}
    &=
    \proj_{\Q_j}(W_j^{k+1}+U_j^k),\\
    U_j^{k+1}
    &=
    U_j^k+W_j^{k+1}-Z_j^{k+1}.
\end{align}
Each block has its own online variable
\[
    u_{j,k}=\log\rho_{j,k},
    \qquad
    u_k=(u_{1,k},\ldots,u_{m,k})\in U^m.
\]
The blockwise residual loss is
\begin{equation}
    \ell_{j,k}^{\mathrm{res}}(u_j)
    =
    \frac{1}{2}
    \left(
        u_j+\log(\norm{Z_j^{k+1}-Z_j^k}+\epsilon)
        -\log(\norm{W_j^{k+1}-Z_j^{k+1}}+\epsilon)
    \right)^2.
\end{equation}
The vector update is coordinatewise projected OGD:
\begin{equation}
    u_{j,k+1}
    =
    \Pi_{[u_{\min},u_{\max}]}
    \left(u_{j,k}-\eta_k \widetilde g_{j,k}\right),
    \qquad j=1,\ldots,m.
    \label{eq:vector-rho-update}
\end{equation}
After changing \(\rho_{j,k}\), the scaled dual must be rescaled per block:
\begin{equation}
    U_j^{k+1}
    \leftarrow
    \frac{\rho_{j,k}}{\rho_{j,k+1}}U_j^{k+1}.
    \label{eq:vector-dual-rescale}
\end{equation}

For the LLM-style PTQ quadratic block, the vector-rho closed form is
\begin{equation}
    \left(
        \frac{X_j^\top X_j}{n_j}
        +(\gamma+\rho_{j,k})I
    \right)
    W_j^{k+1}
    =
    \frac{X_j^\top Y_j}{n_j}
    +\gamma W_j^{\mathrm{fp}}
    +\rho_{j,k}(Z_j^k-U_j^k).
    \label{eq:vector-llm-ptq-w-update}
\end{equation}
This is the same argmin as the scalar case, but with a block-specific penalty.
The proof extension is straightforward at the perturbation layer:
\begin{equation}
    \norm{u_{k+1}-u_k}_2
    \le
    \eta_k \sqrt{m}\,G_u
\end{equation}
when every coordinate gradient is clipped by \(G_u\).  The possible downside is
statistical and algorithmic: the vector method has more degrees of freedom,
so it needs reliable blockwise task signals or curvature normalization to avoid
chasing noisy per-block metrics.

\section{Quantization Split}

The quantization formulation separates continuous trainable weights from a
deployable quantized copy:
\begin{equation}
    \min_{W,Z}\; F(W)+I_{\Q}(Z)
    \quad \text{s.t.}\quad W-Z=0.
    \label{eq:quant-split}
\end{equation}
Here \(F\) is a calibration, fine-tuning, or training objective and
\(I_{\Q}\) is the indicator of the quantization set \(\Q\).  The scaled ADMM
updates are
\begin{align}
    W^{k+1}
    &\approx
    \argmin_W
    F(W)+\frac{\rho_k}{2}\norm{W-Z^k+U^k}^2,\\
    Z^{k+1}
    &=
    \proj_{\Q}(W^{k+1}+U^k),\\
    U^{k+1}
    &=U^k+W^{k+1}-Z^{k+1}.
\end{align}
The projection onto \(\Q\) is nonconvex for low-bit quantization.  For
symmetric \(b\)-bit uniform per-channel quantization,
\begin{equation}
    \Q_b
    =
    \{s_j q_{ij}: q_{ij}\in\{-2^{b-1}+1,\ldots,2^{b-1}-1\},
    \ s_j>0\},
\end{equation}
the implementation uses the standard per-output-channel projection heuristic
\begin{equation}
    s_j=\frac{\max_i |V_{ij}|}{2^{b-1}-1},
    \qquad
    Z_{ij}=s_j\,
    \operatorname{clip}\!\left(
    \operatorname{round}(V_{ij}/s_j),-2^{b-1}+1,2^{b-1}-1
    \right),
    \label{eq:uniform-quant-proj}
\end{equation}
where \(V=W^{k+1}+U^k\).  This is the deployable low-bit projection used in the
experiments; it is a heuristic projection when the scale is also optimized.

\section{Example 1: Nonconvex Quantized Tanh Network}

The first nonconvex example is a one-hidden-layer regression model
\begin{equation}
    \widehat y(W_1,W_2;x)=\tanh(xW_1)W_2.
\end{equation}
The ADMM split is
\begin{equation}
    \min_{W_1,W_2,Z_1,Z_2}
    \frac{1}{2n}\norm{\tanh(XW_1)W_2-Y}^2
    +\frac{\gamma}{2}(\norm{W_1}^2+\norm{W_2}^2)
    +I_{\Q_b}(Z_1)+I_{\Q_b}(Z_2)
\end{equation}
\begin{equation}
    \text{s.t.}\qquad W_1-Z_1=0,\qquad W_2-Z_2=0.
\end{equation}
The continuous \(W\)-update is nonconvex, so it is solved inexactly by several
gradient steps on
\begin{align}
    \Phi_k(W_1,W_2)
    &=
    \frac{1}{2n}\norm{\tanh(XW_1)W_2-Y}^2
    +\frac{\gamma}{2}(\norm{W_1}^2+\norm{W_2}^2)\\
    &\quad
    +\frac{\rho_k}{2}\norm{W_1-Z_1^k+U_1^k}^2
    +\frac{\rho_k}{2}\norm{W_2-Z_2^k+U_2^k}^2.
\end{align}
Then
\begin{equation}
    Z_i^{k+1}=\proj_{\Q_b}(W_i^{k+1}+U_i^k),
    \qquad
    U_i^{k+1}=U_i^k+W_i^{k+1}-Z_i^{k+1},
    \qquad i=1,2.
\end{equation}
This is the closest small-scale analogue of quantized nonconvex training in
the current codebase.

\section{Example 2: Synthetic LLM-Style Blockwise PTQ}

The second example models layerwise post-training quantization for transformer
linear blocks.  For each block \(j\), calibration activations \(X_j\) and a
full-precision weight \(W_j^{\mathrm{fp}}\) define targets
\begin{equation}
    Y_j=X_jW_j^{\mathrm{fp}}.
\end{equation}
The split problem is
\begin{equation}
    \min_{\{W_j,Z_j\}}
    \sum_j
    \frac{1}{2n_j}\norm{X_jW_j-Y_j}^2
    +\frac{\gamma}{2}\norm{W_j-W_j^{\mathrm{fp}}}^2
    +I_{\Q_b}(Z_j)
    \quad\text{s.t.}\quad
    W_j-Z_j=0.
    \label{eq:llm-ptq}
\end{equation}
This is nonconvex because \(Z_j\in\Q_b\), but the continuous \(W_j\)-update is
quadratic and has the closed form
\begin{equation}
    \left(
        \frac{X_j^\top X_j}{n_j}
        +(\gamma+\rho_k)I
    \right)
    W_j^{k+1}
    =
    \frac{X_j^\top Y_j}{n_j}
    +\gamma W_j^{\mathrm{fp}}
    +\rho_k(Z_j^k-U_j^k).
    \label{eq:llm-ptq-w-update}
\end{equation}
The quantized block update is
\begin{equation}
    Z_j^{k+1}=\proj_{\Q_b}(W_j^{k+1}+U_j^k),
    \qquad
    U_j^{k+1}=U_j^k+W_j^{k+1}-Z_j^{k+1}.
\end{equation}
The synthetic blocks mimic transformer projections:
\[
\texttt{q,k,v,o\_proj},\quad
\texttt{gate\_proj},\quad
\texttt{up\_proj},\quad
\texttt{down\_proj}.
\]
The calibration activations are heavy-tailed to mimic outlier channels, and the
weights contain rare large entries to stress the quantizer.

\section{Experimental Results and Conclusions}

All reported results in this section use three seeds, 100 ADMM iterations, and
4-bit symmetric per-channel quantization.  The deploy metric for the LLM-style
PTQ example is relative calibration error; lower is better.

\subsection{Robustness to Initial \texorpdfstring{\(\rho\)}{rho}}

The initial-penalty sweep used
\(\rho_0\in\{10^{-2},10^{-1},1,10,10^2\}\).  Table
\ref{tab:rho-robustness} reports mean deploy relative error across this grid,
as well as the best and worst grid values.

\begin{table}[h]
\centering
\caption{Initialization robustness on synthetic LLM PTQ.}
\label{tab:rho-robustness}
\begin{tabular}{lcccc}
\toprule
Method & Mean deploy & Best deploy & Worst deploy & Relative range\\
\midrule
\texttt{online\_ogd} & 0.352 & 0.349 & 0.354 & 1.4\%\\
\texttt{online\_ogd\_task\_aware} & 0.343 & 0.339 & 0.347 & 2.3\%\\
\texttt{online\_ogd\_task\_feasibility} & 0.323 & 0.316 & 0.332 & 5.1\%\\
\texttt{vector\_online\_ogd} & 0.350 & 0.347 & 0.354 & 2.0\%\\
\texttt{vector\_task\_feasibility} & 0.324 & 0.317 & 0.332 & 4.4\%\\
\bottomrule
\end{tabular}
\end{table}

The scalar online methods are fairly robust with respect to initialization:
they usually recover similar final penalties even from poor starts.  However,
the most accurate method is not the most initialization-insensitive method.
This separates two claims:
\[
    \text{robust penalty path}
    \;\not\Rightarrow\;
    \text{best deploy accuracy}.
\]
The vector-rho implementation is principled, but it only modestly changes the
result in the present synthetic PTQ setting.  Its main value is diagnostic:
it can reveal block heterogeneity, but it needs better blockwise task signals
or curvature normalization before it becomes a clear accuracy win.

\subsection{Alternative Online Losses for LLM PTQ}

Table \ref{tab:llm-loss-sweep} compares several scalar online losses.  The
new variants include:
\begin{itemize}
    \item \texttt{sum\_log\_residuals}: minimize
    \(\log\norm{r^{k+1}}+\log\norm{s^{k+1}}\), where
    \(s^{k+1}=\rho_k(Z^{k+1}-Z^k)\).
    \item \texttt{norm\_magnitude}: minimize squared normalized log residual
    magnitudes.
    \item \texttt{task\_feas}: combine the task secant loss with stronger
    primal-feasibility pressure.
\end{itemize}

\begin{table}[h]
\centering
\caption{Alternative online losses on synthetic LLM PTQ.}
\label{tab:llm-loss-sweep}
\begin{tabular}{lcccc}
\toprule
Method & Deploy error & Final primal & Final dual & Final \(\rho\)\\
\midrule
\texttt{fixed\_rho\_10} & 0.278 & 2.19 & 28.21 & 10.0\\
\texttt{online\_ogd\_balance} & 0.349 & 7.64 & 7.75 & 0.72\\
\texttt{online\_ogd\_task\_feasibility} & 0.316 & 5.79 & 15.80 & 2.05\\
\texttt{online\_ogd\_task\_feas\_mid} & 0.282 & 2.97 & 23.72 & 6.35\\
\texttt{online\_ogd\_task\_feas\_aggressive} & \textbf{0.254} & \textbf{1.07} & 37.37 & 29.42\\
\texttt{online\_ogd\_sum\_log\_residuals} & 0.364 & 9.10 & 1.54 & 0.109\\
\texttt{online\_ogd\_norm\_magnitude} & 0.368 & 8.72 & 0.041 & 0.0027\\
\bottomrule
\end{tabular}
\end{table}

The simple sum-log residual loss is not competitive.  The reason is visible
from the surrogate derivative.  With current residuals treated as fixed,
\[
    \ell_k(u)
    =
    \log(\norm{r^{k+1}}+\epsilon)
    +
    \log(e^u\norm{d^{k+1}}+\epsilon)
\]
has approximately
\[
    \frac{\partial \ell_k}{\partial u}\approx 1
\]
whenever the dual residual is nonzero.  Projected OGD therefore shrinks
\(\rho\) almost monotonically.  In the PTQ experiment this weakens the
continuous/quantized coupling and hurts deploy accuracy.

\subsection{Accuracy--Convergence Tradeoff}

The LLM PTQ results show a real tradeoff.  Residual-balanced OGD has a smaller
final max ADMM residual than aggressive task-feasibility tuning, but it has
worse deploy accuracy.  Conversely, the aggressive task-feasibility loss
improves deploy error from \(0.278\) for fixed \(\rho=10\) to \(0.254\), and
reduces the primal residual from \(2.19\) to \(1.07\), but increases the dual
residual from \(28.21\) to \(37.37\).

Thus ``faster convergence'' must be specified carefully.  If convergence means
lower primal mismatch \(W-Z\), then stronger online penalties can help and may
be worth it for PTQ because the deployable quantized copy is the object of
interest.  If convergence means a balanced KKT-like ADMM residual, then the
aggressive rule is not better.  For quantization, the practical objective is
often deploy accuracy under a calibration budget, so a Pareto view is more
appropriate than choosing the method with the smallest residual alone.

\subsection{Comparison to LLM Quantization Proxy Baselines}

To check whether the ADMM quantizer is competitive with existing LLM PTQ
ideas, the experiments also include local proxy baselines inspired by common
methods:
\begin{itemize}
    \item \textbf{RTN}: per-output-channel round-to-nearest uniform
    quantization.
    \item \textbf{GPTQ-like}: sequential Hessian-compensated quantization using
    the calibration Gram matrix \(X^\top X/n\).
    \item \textbf{AWQ/SmoothQuant-like}: activation-aware input-channel
    rescaling grids before uniform quantization.
    \item \textbf{Hessian diagonal clipping}: scale selection by a diagonal
    Hessian-weighted reconstruction objective.
\end{itemize}
These are compact local implementations on the synthetic calibration problem,
not official library implementations.  They are nevertheless useful because
they test the core issue: whether pure ADMM with uniform projection can compete
against curvature-aware PTQ.

\begin{table}[h]
\centering
\caption{ADMM quantization versus local LLM PTQ proxy baselines.}
\label{tab:llm-proxy-baselines}
\begin{tabular}{lccc}
\toprule
Method & Family & Deploy error & Weight error\\
\midrule
\texttt{gptq\_like\_sequential} & Hessian PTQ & \textbf{0.205} & 0.343\\
\texttt{hessian\_diag\_clip} & Hessian PTQ & 0.222 & \textbf{0.229}\\
\texttt{rtn\_per\_channel} & RTN & 0.252 & 0.255\\
\texttt{admm\_task\_feas\_aggressive} & ADMM & 0.254 & 0.266\\
\texttt{admm\_fixed\_rho\_10} & ADMM & 0.278 & 0.264\\
\texttt{admm\_task\_feasibility} & ADMM & 0.316 & 0.286\\
\bottomrule
\end{tabular}
\end{table}

This result is important.  The current ADMM method, even with aggressive online
penalty tuning, does not beat Hessian-aware PTQ.  It is roughly competitive
with RTN, but GPTQ-like error compensation is much stronger on the calibration
objective.  Therefore the right conclusion is not that online ADMM is already a
state-of-the-art LLM quantizer.  The right conclusion is:
\[
    \text{online penalty tuning helps ADMM,}
    \qquad
    \text{but the } Z\text{-projection must become curvature-aware.}
\]
The most promising next algorithm is an ADMM framework whose \(Z\)-update uses
a GPTQ/AWQ-style calibrated quantization step, while the online controller tunes
the coupling strength and feasibility/accuracy tradeoff.

\subsection{Curvature-Aware ADMM Follow-Up}

The proxy baseline result motivates replacing the Euclidean quantized
projection by a curvature-aware quantized \(Z\)-update.  For block \(j\), let
\[
    H_j=\frac{X_j^\top X_j}{n_j},
    \qquad
    V_j^k=W_j^{k+1}+U_j^k.
\]
Instead of projecting \(V_j^k\) onto \(\Q_b\) in Euclidean norm, use the
calibrated quadratic surrogate
\begin{equation}
    Z_j^{k+1}
    \approx
    \argmin_{Z_j\in\Q_b}
    \frac{\alpha}{2}
    \norm{X_j(Z_j-W_j^{\mathrm{fp}})}_F^2
    +
    \frac{\rho_k}{2}
    \norm{Z_j-V_j^k}_F^2.
    \label{eq:curvature-aware-z}
\end{equation}
The continuous minimizer of \eqref{eq:curvature-aware-z} is
\begin{equation}
    T_j^k
    =
    (\alpha H_j+\rho_k I)^{-1}
    (\alpha H_j W_j^{\mathrm{fp}}+\rho_k V_j^k).
    \label{eq:curvature-aware-target}
\end{equation}
The implementation then quantizes \(T_j^k\) using one of three approximations:
diagonal Hessian clipping, GPTQ-like sequential error compensation, or
activation-aware scaling.  This changes the method from exact ADMM for the
original split into a majorized/curvature-aware ADMM heuristic.  That is an
acceptable experimental direction, but the theorem should state the extra
projection error explicitly.

\begin{table}[h]
\centering
\caption{Curvature-aware ADMM on synthetic LLM PTQ.}
\label{tab:curvature-admm}
\begin{tabular}{lccc}
\toprule
Method & Deploy error & Final primal & Final \(\rho\)\\
\midrule
\texttt{gptq\_like\_sequential} & \textbf{0.205} & -- & --\\
\texttt{admm\_gptq\_like\_online\_aggressive} & 0.214 & \textbf{0.494} & 26.82\\
\texttt{hessian\_diag\_clip} & 0.222 & -- & --\\
\texttt{admm\_hessian\_diag\_fixed10} & 0.236 & 1.19 & 10.0\\
\texttt{admm\_hessian\_diag\_online\_aggressive} & 0.237 & 0.640 & 25.88\\
\texttt{admm\_uniform\_online\_aggressive} & 0.254 & 1.07 & 29.42\\
\texttt{rtn\_per\_channel} & 0.252 & -- & --\\
\bottomrule
\end{tabular}
\end{table}

This is the strongest evidence so far for the method.  Plain ADMM with uniform
projection is only RTN-level.  Curvature-aware ADMM with online penalty tuning
nearly reaches the standalone GPTQ-like proxy while maintaining the ADMM
structure and producing a much smaller primal residual.  The research message
should therefore be:
\[
    \text{online ADMM is promising when coupled to a strong calibrated}
    \ Z\text{-update.}
\]
The remaining gap to GPTQ-like PTQ is small enough to justify a real-model
perplexity experiment.

\subsection{Next Online Losses to Test}

The next experiments should test online losses that explicitly encode this
Pareto structure:
\begin{enumerate}
    \item \textbf{Constrained task loss.}  Minimize the task secant loss while
    penalizing only residual violations above a budget:
    \[
        \ell_k(u)=\log(m_k+\epsilon)
        +\alpha [\log(\norm{r^{k+1}}/\varepsilon_{\rm pri})]_+^2
        +\beta [\log(\norm{s^{k+1}}/\varepsilon_{\rm dual})-\tau]_+^2.
    \]
    \item \textbf{Two-sided bandit probing.}  Occasionally evaluate short
    copied-state rollouts at \(u_k+\delta\) and \(u_k-\delta\), then estimate
    a deploy-loss gradient by finite differences.  This is more expensive but
    less biased than the one-step secant.
    \item \textbf{Curvature-normalized block loss.}  For vector rho, normalize
    each block residual by a curvature proxy such as
    \(\operatorname{tr}(X_j^\top X_j/n_j)\) or the spectral norm of the
    calibration Gram matrix.
    \item \textbf{Dual-budget task feasibility.}  Keep the aggressive
    primal-feasibility pressure, but add a soft cap on dual residual growth so
    the method searches along the deploy-accuracy/residual Pareto frontier
    instead of simply increasing \(\rho\).
\end{enumerate}

\section{Theorems to Prove}

\subsection{Experiment-to-Theorem Map}

The experiments now correspond to distinct theorem components rather than one
single monolithic claim:
\begin{table}[h]
\centering
\caption{Algorithmic modifications and corresponding theorem targets.}
\label{tab:theorem-map}
\begin{tabular}{p{0.31\linewidth}p{0.31\linewidth}p{0.28\linewidth}}
\toprule
Modification & Empirical role & Theorem target\\
\midrule
Projected OGD on \(u=\log\rho\) & Makes penalty tuning robust to scale and
initialization & Bounded path length and scalar online-regret perturbation\\
Task-aware secant loss & Aligns penalty updates with deploy metric & Dynamic
regret or linearized online-loss bound\\
Feasibility-task pressure & Improves PTQ deploy accuracy by increasing
coupling & Pareto/stationarity bound with residual budget\\
Vector rho & Allows blockwise penalties for transformer layers & Product-space
path-length bound in \(U^m\)\\
Curvature-aware \(Z\)-update & Makes ADMM competitive with Hessian PTQ & ADMM
descent with inexact/majorized quantized projection error\\
\bottomrule
\end{tabular}
\end{table}

The final theory should not claim global optimality for nonconvex quantization.
A realistic theorem stack is:
\[
    \text{bounded online penalty path}
    \;+\;
    \text{descent up to quantized projection error}
    \;+\;
    \text{stationarity/noise floor}
    \;+\;
    \text{task/residual Pareto interpretation}.
\]

\subsection{Theorem A: Deterministic Nonconvex Proximal ADMM Descent}

\textbf{Assumptions.}
\begin{itemize}
    \item \(f\) is lower bounded and \(L_f\)-smooth on a bounded level set.
    \item \(g\) is proper lower-semicontinuous and either weakly convex,
          prox-regular, or an indicator of a prox-regular set on the relevant
          level set.
    \item The subproblems are solved exactly or with summable/relative errors.
    \item \(\rho_k\in[\rho_{\min},\rho_{\max}]\), and the penalty path has
          bounded cumulative variation \(\sum_k |u_{k+1}-u_k|<\infty\), or a
          controlled finite-horizon bound using \eqref{eq:path-length}.
    \item Proximal parameters \(\tau_x,\tau_z\) dominate the negative
          curvature of the nonconvex blocks.
\end{itemize}

\textbf{Merit function.}  Prove descent for
\begin{equation}
    \Psi_k
    =
    \Lagr_{\rho_k}(x^k,z^k,\lambda^k)
    +a_x\norm{x^k-x^{k-1}}^2
    +a_z\norm{z^k-z^{k-1}}^2
    +a_\lambda\norm{\lambda^k-\lambda^{k-1}}^2.
\end{equation}

\textbf{Target statement.}  Show that
\begin{align}
    \Psi_{k+1}
    &\le
    \Psi_k
    -m_x\norm{x^{k+1}-x^k}^2
    -m_z\norm{z^{k+1}-z^k}^2
    -m_r\norm{r^{k+1}}^2\\
    &\quad
    +c_u|u_{k+1}-u_k|
    +c_e\norm{e_k}^2.
    \label{eq:nonconvex-descent-target}
\end{align}
This replaces the convex objective-gap sandwich.  Summing
\eqref{eq:nonconvex-descent-target} yields best-iterate stationarity and
residual decay up to the penalty-path and inexactness terms.

\subsection{Theorem B: Online Penalty Perturbation Bound}

Assume the online update satisfies \eqref{eq:path-length}.  Since \(e^u\) and
\(e^{-u}\) are Lipschitz on compact \(U\), all terms produced by changing
\(\rho_k\) can be bounded by
\begin{equation}
    \sum_{k=0}^{K-1} C |u_{k+1}-u_k|
    \le
    C G_u\sum_{k=0}^{K-1}\eta_k.
\end{equation}
For \(\eta_k=\eta_0/\sqrt{k+1}\), this gives an \(O(\sqrt K)\) perturbation in
the summed descent inequality and an \(O(K^{-1/2})\) perturbation after
averaging.  For full-sequence nonconvex convergence, strengthen this to
summable variation, finite accepted penalty changes, or an accept/reject
safeguard.

\subsection{Theorem B2: Task-Aware Online Loss Regret}

The task-aware variants require a separate scalar online-learning statement.
Let
\[
    \widetilde g_k
    =
    \clip\left(
        \alpha g_k^{\mathrm{res}}
        +\beta g_k^{\mathrm{task}}
        -\chi s_k^{\mathrm{feas}},
        G_u
    \right)
\]
with \(\chi=0\) for the composite task-secant method.  If
\(\widetilde g_k\) is bounded and \(U=[u_{\min},u_{\max}]\) is compact, then
projected OGD gives the deterministic regret bound
\begin{equation}
    \sum_{k=0}^{K-1}
    \ip{\widetilde g_k}{u_k-u}
    \le
    \frac{\diam(U)^2}{2\eta_{K-1}}
    +
    \frac{1}{2}\sum_{k=0}^{K-1}\eta_k G_u^2,
    \qquad \forall u\in U.
    \label{eq:task-aware-regret}
\end{equation}
This theorem does not say that the task metric is convex in \(u\).  Instead,
it states that the implemented update is stable and has bounded regret against
the linearized online losses generated by the observed trajectory.  To connect
it to deploy loss, one must add either a local smoothness assumption for
\(\log m(u)\), a dynamic-regret comparator sequence, or an empirical
validation claim.

\subsection{Theorem B3: Curvature-Aware Quantized Projection Error}

For the curvature-aware \(Z\)-update, define the ideal continuous target
\(T_j^k\) by \eqref{eq:curvature-aware-target} and the quantized output
\(Z_j^{k+1}\in\Q_b\).  The proof should assume or establish a projection-error
condition of the form
\begin{equation}
    \sum_j
    \norm{Z_j^{k+1}-T_j^k}_{\alpha H_j+\rho_k I}^2
    \le
    \varepsilon_{q,k}^2
    +
    c_q
    \inf_{Z\in\Q_b}
    \sum_j
    \norm{Z_j-T_j^k}_{\alpha H_j+\rho_k I}^2.
    \label{eq:curvature-projection-error}
\end{equation}
For exact projection, \(c_q=1\) and \(\varepsilon_{q,k}=0\).  For GPTQ-like,
AWQ-like, or diagonal-Hessian approximate projections, \(\varepsilon_{q,k}\)
captures sequential quantization error, scale-grid suboptimality, and
activation-statistic mismatch.  This condition is the bridge between modern
PTQ baselines and the nonconvex ADMM descent theorem.

\subsection{Theorem C: Expected Stationarity With Quantization Noise}

For stochastic or randomized quantization, assume
\begin{equation}
    \mathbb{E}[e_k\mid\mathcal F_k]=0,
    \qquad
    \mathbb{E}\norm{e_k}^2\le \sigma_q^2+\delta_q^2\norm{v_k}^2.
\end{equation}
Then prove a bound of the form
\begin{equation}
    \frac{1}{K}\sum_{k=0}^{K-1}
    \mathbb{E}
    \left[
        \dist(0,\partial\Psi_{k+1})^2+\norm{r^{k+1}}^2
    \right]
    \le
    O(K^{-1/2})+O(\sigma_q^2+\delta_q^2).
    \label{eq:noise-floor}
\end{equation}
This is the theorem that matches stochastic rounding, low-bit optimizer states,
and direct quantized training.

\subsection{Theorem D: Local PL or Quadratic-Growth Fast Rate}

Near a pretrained checkpoint or a good PTQ initialization, assume a local PL,
quadratic-growth, error-bound, or strong-variational-sufficiency condition:
\begin{equation}
    \dist(0,\partial\Psi_k)^2
    \ge
    \mu(\Psi_k-\Psi^\star)
\end{equation}
on the relevant level set.  Then upgrade the expected stationarity result to
a linear-to-floor statement:
\begin{equation}
    \mathbb{E}[\Psi_k-\Psi^\star]
    \le
    (1-\mu')^k(\Psi_0-\Psi^\star)
    +O(\sigma_q^2+\delta_q^2+\varepsilon_q^2).
\end{equation}
This is the theorem most aligned with layerwise PTQ, local fine-tuning, and
curvature-aware quantization methods.

\section{Proof Obligations Checklist}

\begin{enumerate}
    \item Replace convex saddle-point gap lemmas by sufficient descent of
          \(\Psi_k\).
    \item Replace monotonicity of \(\partial g\) by weak monotonicity,
          prox-regularity, or an explicit projection-error inequality.
    \item Derive bounded iterates from bounded level sets of \(\Psi_k\), weight
          decay, or trust-region/proximal terms.
    \item Keep the online-penalty result as a path-length perturbation theorem,
          not as a convex-regret theorem unless the scalar meta-loss assumptions
          are truly satisfied.
    \item For quantization, state whether errors are deterministic bounded
          errors, unbiased stochastic-rounding errors, or relative
          floating-point errors.
    \item Add a local PL/QG theorem only after the general stationarity theorem
          is established.
\end{enumerate}

\end{document}
